%\part*{Delft}
%\addcontentsline{toc}{part}{Delft Partition}

%\setcounter{chapter}{0}
%\renewcommand{\thechapter}{\arabic{chapter} }
%\def\chaptername{Chapter}

\chapter{Introduction}

This chapter introduces all aspects that form the context and
motivations of the graduate project. We will describe \ldots

\section{Peer-to-Peer architectures}

\subsection{Drawbacks of client-server model}
The traditional client-server computing model is not adequate to address reliability and scalability properties. Even when servers are replicated for fault-tolerance, the synchronization mechanism needed to preserve replica consistency is a major source of overhead.

\subsection{P2P approach}
The peer-to-peer computing model represents a radically different and a plausible alternative for many large scale applications in internet-wide settings. In this model resources are shared between end-user processes themselves in a peer style, meaning that every such process potentially acts both as client and a server.

\subsubsection{Advantages}
Central point of failures disappear as well as the associated performance bottlenecks.

Scalability is achieved because the load is balanced between all processes of the system.

\subsubsection{Disvantages}
In addition, each process requires only local knowledge of the state of the system.

\section{Epidemic protocols}
Epidemic algorithms have been recently recognized as robust and scalable means to disseminate information in large-scale settings. Information is disseminated reliably in a distributed system the same way an epidemic would be propagated throughout a group of individuals: each process of the system chooses random peers to whom it relays the information it has received.

Epidemic algorithms have recently gained popularity as a robust and scalable way of propagating information in distributed systems.
A process that wishes to disseminate a new piece of information to the system does not send it to a server, or a cluster of servers, in charge of forwarding it, but rather to a set of other peer processes, chosen at random. In turn each of each processes does the same, and also forward the information to randomly selected processes, and so fort

The underlying peer-to-peer communication paradigm is the key to the scalability of the dissemination scheme.
Epidemic algorithms have been studied theoretically and their analysis is built on sound mathematical foundations.

